% Options for packages loaded elsewhere
\PassOptionsToPackage{unicode}{hyperref}
\PassOptionsToPackage{hyphens}{url}
%
\documentclass[
  lualatex,ja=standard]{bxjsarticle}
\usepackage{amsmath,amssymb}
\usepackage{lmodern}
\usepackage{iftex}
\ifPDFTeX
  \usepackage[T1]{fontenc}
  \usepackage[utf8]{inputenc}
  \usepackage{textcomp} % provide euro and other symbols
\else % if luatex or xetex
  \usepackage{unicode-math}
  \defaultfontfeatures{Scale=MatchLowercase}
  \defaultfontfeatures[\rmfamily]{Ligatures=TeX,Scale=1}
\fi
% Use upquote if available, for straight quotes in verbatim environments
\IfFileExists{upquote.sty}{\usepackage{upquote}}{}
\IfFileExists{microtype.sty}{% use microtype if available
  \usepackage[]{microtype}
  \UseMicrotypeSet[protrusion]{basicmath} % disable protrusion for tt fonts
}{}
\makeatletter
\@ifundefined{KOMAClassName}{% if non-KOMA class
  \IfFileExists{parskip.sty}{%
    \usepackage{parskip}
  }{% else
    \setlength{\parindent}{0pt}
    \setlength{\parskip}{6pt plus 2pt minus 1pt}}
}{% if KOMA class
  \KOMAoptions{parskip=half}}
\makeatother
\usepackage{xcolor}
\IfFileExists{xurl.sty}{\usepackage{xurl}}{} % add URL line breaks if available
\IfFileExists{bookmark.sty}{\usepackage{bookmark}}{\usepackage{hyperref}}
\hypersetup{
  hidelinks,
  pdfcreator={LaTeX via pandoc}}
\urlstyle{same} % disable monospaced font for URLs
\usepackage{longtable,booktabs,array}
\usepackage{calc} % for calculating minipage widths
% Correct order of tables after \paragraph or \subparagraph
\usepackage{etoolbox}
\makeatletter
\patchcmd\longtable{\par}{\if@noskipsec\mbox{}\fi\par}{}{}
\makeatother
% Allow footnotes in longtable head/foot
\IfFileExists{footnotehyper.sty}{\usepackage{footnotehyper}}{\usepackage{footnote}}
\makesavenoteenv{longtable}
\setlength{\emergencystretch}{3em} % prevent overfull lines
\providecommand{\tightlist}{%
  \setlength{\itemsep}{0pt}\setlength{\parskip}{0pt}}
\setcounter{secnumdepth}{-\maxdimen} % remove section numbering
\ifLuaTeX
  \usepackage{selnolig}  % disable illegal ligatures
\fi

\author{}
\date{}

\begin{document}

\hypertarget{ux4f0aux6577ux4e88ux60f3-5.1-ux306eux8a3cux660e}{%
\section{伊敷予想 5.1
の証明}\label{ux4f0aux6577ux4e88ux60f3-5.1-ux306eux8a3cux660e}}

本報告書では，Ishiki {[}1{]} の Conjecture 5.1，すなわち濃度
\(\aleph_1\) の離散位相空間 \(D\) に対して \(\operatorname{Met}(D)\)
が完備距離化可能でないことを，ZFC の定理として証明する．

\hypertarget{ux8a18ux53f7ux3068ux8a2dux5b9a}{%
\subsection{1. 記号と設定}\label{ux8a18ux53f7ux3068ux8a2dux5b9a}}

\begin{longtable}[]{@{}
  >{\raggedright\arraybackslash}p{(\columnwidth - 2\tabcolsep) * \real{0.50}}
  >{\raggedright\arraybackslash}p{(\columnwidth - 2\tabcolsep) * \real{0.50}}@{}}
\toprule
記号 & 意味 \\
\midrule
\endhead
\(D\) & 濃度 \(\aleph_1\) の離散位相空間 \\
\(\operatorname{Met}(D)\) & \(D\)
の離散位相を誘導する距離の全体．上限距離位相を備える \\
\(\rho(d_1, d_2)\) &
\(\sup_{x,y\in D} \lvert d_1(x,y) - d_2(x,y)\rvert\)（\(\operatorname{Met}(D)\)
上の上限距離） \\
\([0,1]\) & 実数の閉単位区間 \\
\(\omega_1\) & 最小の非可算順序数 \\
\(\aleph_1\) & 最小の非可算基数 \\
\(\mathfrak{c}\) & 連続体濃度 \(2^{\aleph_0}\) \\
\(S^\complement\) & 部分集合 \(S \subseteq [0,1]\) の補集合
\([0,1] \setminus S\)．本報告書で補集合は常に全体集合 \([0,1]\)
でとる \\
\bottomrule
\end{longtable}

\(\operatorname{Met}(D)\) は集合 \(D\)
上の距離すべての集合ではない．離散でない位相を誘導する距離は
\(\operatorname{Met}(D)\) に属さない．

基数については，Cantor の定理より \(\aleph_0 < \mathfrak{c}\)
であり，\(\aleph_1\) は最小の非可算基数だから
\(\aleph_1 \le \mathfrak{c}\) が ZFC で成り立つ．

\textbf{定理 1.1（完備距離化可能部分空間の \(G_\delta\) 性 {[}2, Theorem
3.11{]}）.} \(X\) を距離化可能空間，\(S \subseteq X\)
を相対位相に関して完備距離化可能な部分空間とする．このとき \(S\) は
\(X\) の \(G_\delta\) 集合である．

本報告書では，この定理を対偶の形で \(X = [0,1]\) に対してのみ用いる（§4
ステップ 3）．同じ一般形は {[}3, Theorems 4.3.23--4.3.24{]} および {[}4,
Theorem 24.12{]} にもある．

\textbf{定理 1.2（閉部分空間の完備距離化可能性）.} \(X\)
を完備距離化可能な空間，\(Y \subseteq X\) を閉部分集合とする．このとき
\(Y\) は相対位相に関して完備距離化可能である．

\emph{証明の方針.} \(X\) 上の完備距離の制限が閉部分空間 \(Y\)
上で完備になることから従う．

本報告書で \(\operatorname{Met}(D)\)
に適用する完備距離化可能性についての事実は定理 1.2 のみである（§4
ステップ 5）．

\textbf{Remark（上限距離の値域）.} 上限距離 \(\rho(d_1, d_2)\)
は，\(d_1, d_2\) が非有界で差が非有界のとき \(+\infty\)
をとりうる．しかし \(\min(\rho, 1)\) は \(\rho\)
と同一の位相を定める有限値の距離だから，\(\operatorname{Met}(D)\)
は通常の意味で距離化可能であり，完備距離化可能性は位相のみで定まる性質だから，このことは以下の議論に影響しない．なお
§3 で構成する距離族の値域はすべて \([0,2]\) に含まれ，そこで現れる
\(\rho\) の値はすべて有限である．

\hypertarget{ux4e3bux5b9aux7406}{%
\subsection{2. 主定理}\label{ux4e3bux5b9aux7406}}

\textbf{定理 2.1（伊敷予想 5.1）.} \(\lvert D\rvert = \aleph_1\)
ならば，\(\operatorname{Met}(D)\)
は上限距離位相に関して完備距離化可能でない．

証明は §4 で与える．§3 で必要な補題を準備する．

\hypertarget{ux88dcux984c}{%
\subsection{3. 補題}\label{ux88dcux984c}}

\hypertarget{f_sigma-ux3067ux306aux3044ux6fc3ux5ea6-aleph_1-ux306eux96c6ux5408}{%
\subsubsection{\texorpdfstring{3.1 \(F_\sigma\) でない濃度 \(\aleph_1\)
の集合}{3.1 F\_\textbackslash sigma でない濃度 \textbackslash aleph\_1 の集合}}\label{f_sigma-ux3067ux306aux3044ux6fc3ux5ea6-aleph_1-ux306eux96c6ux5408}}

\textbf{補題 3.1（閉集合の濃度）.} \(F \subseteq [0,1]\)
を閉集合とする．\(\lvert F\rvert < \mathfrak{c}\) ならば \(F\)
は可算である．

\emph{証明.} 対偶を示す．\(F\) を非可算な閉集合とする．\(F\) は
\([0,1]\) の閉部分集合だからそれ自身 Polish
空間であり，Cantor--Bendixson の定理 {[}2, Theorem 6.4{]} より完全集合
\(P\) と可算集合 \(C\) の非交和 \(F = P \sqcup C\) に分解できる．\(F\)
が非可算で \(C\) が可算だから \(P \ne \emptyset\) であり，空でない完全
Polish 空間は濃度 \(\mathfrak{c}\) を持つ {[}2, Theorem 6.2, Corollary
6.3{]}．よって \(\lvert F\rvert \ge \lvert P\rvert = \mathfrak{c}\)
である．

\textbf{補題 3.2（\(F_\sigma\) でない集合の存在）.} \(F_\sigma\)
でない濃度 \(\aleph_1\) の集合 \(A \subseteq [0,1]\) が存在する．

\emph{証明.} 連続体仮説
\(\mathsf{CH}\)（\(\aleph_1 = \mathfrak{c}\)）の成否で場合を分ける．

\textbf{場合 1（\(\mathsf{CH}\) が成立するとき）.}
\(A := [0,1] \setminus \mathbf{Q}\)（\([0,1]\)
内の無理数全体）とおく．\(\lvert[0,1]\rvert = \mathfrak{c}\) で
\(\mathbf{Q} \cap [0,1]\) は可算だから \(\lvert A\rvert = \mathfrak{c}\)
であり，\(\mathsf{CH}\) より \(\lvert A\rvert = \aleph_1\)．

\(A\) が \(F_\sigma\) でないことを示す（この部分は \(\mathsf{CH}\)
を用いない）．\(A = \bigcup_{n\in\mathbf{N}} F_n\)（各 \(F_n\) は
\([0,1]\) の閉集合）と仮定する．\(\mathbf{Q} \cap [0,1]\) は \([0,1]\)
で稠密で \(F_n \cap \mathbf{Q} = \emptyset\) だから，各 \(F_n\)
は内点を持たない閉集合，したがって疎（nowhere
dense）である．\(\mathbf{Q} \cap [0,1]\)
も可算個の一点集合の合併として痩集合（meager
set）だから，\([0,1] = A \cup (\mathbf{Q} \cap [0,1])\) は \([0,1]\)
の痩集合となり，完備距離空間は自身の中で痩せていないという Baire
のカテゴリ定理 {[}4, Corollary 25.4{]} に反する．よって \(A\) は
\(F_\sigma\) でない．

\textbf{場合 2（\(\mathsf{CH}\)
が成立しないとき，\(\aleph_1 < \mathfrak{c}\)）.}
\(\aleph_1 \le \mathfrak{c} = \lvert[0,1]\rvert\) だから，濃度
\(\aleph_1\) の部分集合 \(A \subseteq [0,1]\)
が存在する．\(A = \bigcup_{n\in\mathbf{N}} F_n\)（各 \(F_n\)
は閉集合）と仮定すると，各 \(n\) について
\(\lvert F_n\rvert \le \lvert A\rvert = \aleph_1 < \mathfrak{c}\)
だから，補題 3.1 より \(F_n\) は可算である．よって \(A\)
は可算個の可算集合の合併として可算となり，\(\lvert A\rvert = \aleph_1\)
に矛盾する．ゆえに \(A\) は \(F_\sigma\) でない．

\textbf{Remark（場合分けを避ける構成）.} \(\mathsf{CH}\)
による場合分けを避けるなら，Bernstein 集合による一様な構成もある．\(B\)
と \(B^\complement\) がともに \([0,1]\)
のすべての非可算閉集合と交わるような集合
\(B \subseteq [0,1]\)（Bernstein
集合，\(\lvert B\rvert = \mathfrak{c}\)）が超限帰納法で構成できる {[}5,
p.~24{]}．濃度 \(\aleph_1\) の部分集合 \(A \subseteq B\)
をとる．\(A = \bigcup_n F_n\)（各 \(F_n\) は閉集合）と書ければ，\(A\)
は非可算だからある \(F_n\) は非可算閉集合であり，\(F_n \subseteq B\) は
\(B^\complement\) がすべての非可算閉集合と交わることに反する．よってこの
\(A\) も \(F_\sigma\) でなく，この議論は \(\mathsf{CH}\)
の成否によらない．

\hypertarget{ux8a2dux5b9aux306eux6b63ux898fux5316}{%
\subsubsection{3.2
設定の正規化}\label{ux8a2dux5b9aux306eux6b63ux898fux5316}}

補題 3.3 以降では，\(D\) を §3.3
で定める具体的な集合にとる．これで一般性を失わないことを先に確認する．離散位相空間の位相型はその濃度のみで定まり，離散空間の間の全単射
\(f : D \to D'\) は同相写像である．対応
\(d \mapsto d \circ (f \times f)\) は \(D'\) 上の距離を \(D\)
上の距離に移し，誘導する位相を保つから，\(\operatorname{Met}(D')\) から
\(\operatorname{Met}(D)\) への全単射を与える（逆写像は \(f^{-1}\)
による同じ対応）．さらに \(\sup\) をとる添字集合が \(f \times f\)
で全単射に移り合うから，この対応は上限距離 \(\rho\)
を保つ．すなわち等長同型である．したがって \(\operatorname{Met}(D)\)
の距離空間としての（特に位相空間としての）同型類は \(\lvert D\rvert\)
のみで定まる．

\hypertarget{ux8dddux96e2ux65cfux306eux69cbux6210}{%
\subsubsection{3.3
距離族の構成}\label{ux8dddux96e2ux65cfux306eux69cbux6210}}

補題 3.2 の与える
\(A \subseteq [0,1]\)（\(\lvert A\rvert = \aleph_1\)，\(A\) は
\(F_\sigma\) でない）と，単射な列挙
\[A = \{a_\alpha : \alpha < \omega_1\}, \qquad \alpha \ne \beta \implies a_\alpha \ne a_\beta\]
を固定する．

\textbf{補題 3.3（距離族の構成）.} 濃度 \(\aleph_1\) の離散空間 \(D\)
と，\(D\) 上の距離の族 \((d_t)_{t \in [0,1]}\) であって
\[d_t \in \operatorname{Met}(D) \iff t \in A^\complement\]
を満たすものが存在する．

\textbf{構成.} 台集合を
\[D := \{p, q\} \ \sqcup\ \{y_\alpha : \alpha < \omega_1\} \ \sqcup\ \{x_{\alpha,n} : \alpha < \omega_1,\ n \in \mathbf{N}\}\]
（三つの部分は互いに交わらず，各添字は相異なる元を与える）で定め，離散位相を与える．濃度は
\[\lvert D\rvert = 2 + \aleph_1 + \aleph_1 \cdot \aleph_0 = \aleph_1\]
である．\(\alpha < \omega_1\) に対して
\[C_\alpha := \{y_\alpha\} \cup \{x_{\alpha,n} : n \in \mathbf{N}\}\]
とおくと， \[D = \{p,q\} \ \sqcup \bigsqcup_{\alpha<\omega_1} C_\alpha\]
は \(D\) の分割である．

\(t \in [0,1]\)，\(\alpha < \omega_1\)，\(n \in \mathbf{N}\) に対して
\[s_n(\alpha, t) := \min\bigl(1,\ \lvert t - a_\alpha\rvert + 1/(n+1)\bigr)\]
とおく．\(0 < s_n(\alpha,t) \le 1\) である．

\(t \in [0,1]\) に対し，写像 \(d_t : D \times D \to [0,2]\)
を，\(d_t(u,u) := 0\) と，相異なる 2
点の非順序対に対する次の表とで定める（表は非順序対に値を与えるから，\(d_t\)
は対称である）．

\begin{longtable}[]{@{}ll@{}}
\toprule
非順序対 \(\{u,v\}\) & \(d_t(u,v)\) \\
\midrule
\endhead
\(\{p,\ q\}\) & \(1 + t\) \\
\(\{y_\alpha,\ x_{\alpha,n}\}\) & \(s_n(\alpha,t)\) \\
\(\{x_{\alpha,n},\ x_{\alpha,m}\}\)（\(n \ne m\)） &
\(s_n(\alpha,t) + s_m(\alpha,t)\) \\
上記以外（分割の相異なる成分に属する 2 点） & \(1\) \\
\bottomrule
\end{longtable}

補題 3.3 の主張を，(C1) \(d_t\) は距離である，(C2)
\(t \in A^\complement\) ならば \(d_t \in \operatorname{Met}(D)\)，(C3)
\(t \in A\) ならば
\(d_t \notin \operatorname{Met}(D)\)，の三つに分けて示す．

\textbf{(C1) 任意の \(t \in [0,1]\) に対し，\(d_t\) は \(D\)
上の距離である．}

\emph{証明.} 対称性と \(d_t(u,u) = 0\)
は定義による．表の各値は正だから，\(d_t(u,v) = 0\) と \(u = v\)
は同値である．

三角不等式 \(d_t(u,w) \le d_t(u,v) + d_t(v,w)\) は，\(u,v,w\)
のうちに一致する 2 点があれば自明だから，相異なる 3
点について示す．\(\alpha\) を固定した文脈では \(s_n := s_n(\alpha,t)\)
と略記する．表より，相異なる 2 点の距離は常に \(2\)
以下であり，分割の相異なる成分に属する 2 点の距離は \(1\) である．3
点が分割とどう交わるかで場合を分ける．

\textbf{場合 (a)　3 点が分割の同一成分に属する.} 3 点が属する成分は 2
元集合 \(\{p, q\}\) ではなく，ある \(C_\alpha\) である．3 点が
\(y_\alpha,\ x_{\alpha,n},\ x_{\alpha,m}\)（\(n \ne m\)）のとき，検証すべき不等式は
\[s_n + s_m \le s_n + s_m \quad (y_\alpha \text{ 経由}), \qquad s_n \le (s_n + s_m) + s_m \quad (x_{\alpha,m} \text{ 経由})\]
の 2 つ（残りはその対称形）であり，いずれも成り立つ．3 点が
\(x_{\alpha,n},\ x_{\alpha,m},\ x_{\alpha,k}\)（添字は相異なる）のときは
\[s_n + s_m \le (s_n + s_k) + (s_k + s_m).\]

\textbf{場合 (b)　ちょうど 2 点が分割の同一成分 \(E\) に属する.}
\(u, v \in E\)，\(z \notin E\) とすると \(d_t(u,z) = d_t(v,z) = 1\)
だから
\[d_t(u,v) \le 2 = d_t(u,z) + d_t(z,v), \qquad d_t(u,z) = 1 \le d_t(u,v) + d_t(v,z).\]

\textbf{場合 (c)　3 点が分割の相異なる 3 成分に属する.} 3
つの対の距離はすべて \(1\) であり，\(1 \le 1 + 1\)．

場合 (a)--(c) は 3 点の配置を尽くすから，三角不等式が成り立つ．

\textbf{(C2) \(t \in A^\complement\) ならば
\(d_t \in \operatorname{Met}(D)\) である．}

\emph{証明.} 距離 \(d\) の誘導する位相が離散位相であることは，各
\(u \in D\) に対して
\[\delta(u) := \inf_{v \in D,\ v \ne u} d(u,v) > 0\]
が成り立つことと同値である．\(t \in A^\complement\)
を固定し，各点でこれを確かめる．

\begin{itemize}
\tightlist
\item
  \(u = p\) のとき，\(d_t(p,q) = 1+t \ge 1\) で他のすべての点への距離は
  \(1\) だから \(\delta(p) = 1\)．\(u = q\) も同様に \(\delta(q) = 1\)．
\item
  \(u = x_{\alpha,n}\) のとき，\(u\) から出る距離の値は
  \(s_n(\alpha,t)\)（\(y_\alpha\)
  へ），\(s_n(\alpha,t) + s_m(\alpha,t)\)（\(x_{\alpha,m}\)，\(m \ne n\)
  へ），\(1\)（\(C_\alpha\) の外へ）だから，\(s_n(\alpha,t) \le 1\) より
  \[\delta(x_{\alpha,n}) = s_n(\alpha,t) \ge 1/(n+1) > 0.\]
\item
  \(u = y_\alpha\) のとき，値は
  \(s_n(\alpha,t)\)（\(n \in \mathbf{N}\)）と \(1\)
  である．\(n \to \infty\) のとき
  \(\lvert t-a_\alpha\rvert + 1/(n+1) \downarrow \lvert t-a_\alpha\rvert\)
  であり，\(r \mapsto \min(1,r)\) は連続かつ単調非減少だから
  \[\delta(y_\alpha) = \min\bigl(1, \lvert t - a_\alpha\rvert\bigr).\]
  \(t \in A^\complement\) かつ \(a_\alpha \in A\) より
  \(t \ne a_\alpha\) だから \(\delta(y_\alpha) > 0\)．
\end{itemize}

よって \(d_t\) の誘導する位相は離散位相であり，(C1) と合わせて
\(d_t \in \operatorname{Met}(D)\)．

\textbf{(C3) \(t \in A\) ならば \(d_t \notin \operatorname{Met}(D)\)
である．}

\emph{証明.} \(t = a_\beta\)（\(\beta < \omega_1\)）とする．(C1) より
\(d_{a_\beta}\) は \(D\) 上の距離である．一方
\[s_n(\beta, a_\beta) = \min\bigl(1,\ 1/(n+1)\bigr) = 1/(n+1)\] だから
\(d_{a_\beta}(y_\beta, x_{\beta,n}) = 1/(n+1)\) であり，任意の
\(\varepsilon > 0\) に対し \(1/(n+1) < \varepsilon\) となる
\(n \in \mathbf{N}\) をとれば，\(y_\beta\) の \(\varepsilon\)-近傍は
\(x_{\beta,n} \ne y_\beta\) を含む．よって \(\{y_\beta\}\)
は開でなく，\(d_{a_\beta}\) の誘導する位相は離散位相と異なる．

(C2) と (C3) より
\(d_t \in \operatorname{Met}(D) \iff t \in A^\complement\)
が成り立ち，補題 3.3 が示された．

\hypertarget{ux9589ux57cbux3081ux8fbcux307f}{%
\subsubsection{3.4 閉埋め込み}\label{ux9589ux57cbux3081ux8fbcux307f}}

\textbf{補題 3.4（閉埋め込み）.} \(A^\complement\) に \([0,1]\)
からの相対位相を与える．写像
\(\varphi : A^\complement \to \operatorname{Met}(D)\)，\(t \mapsto d_t\)
は双リプシッツ埋め込みであり，像 \(\varphi(A^\complement)\) は
\(\operatorname{Met}(D)\) の閉集合である．

補題 3.3 の \(A\)，\(D\)，\((d_t)_{t \in [0,1]}\)
をそのまま用いる．距離族の値域が \([0,2]\) に含まれるから，以下に現れる
\(\rho\) の値はすべて有限である．主張を (C4)（双リプシッツ評価）と
(C5)（像の閉性）に分けて示す．

\textbf{(C4) 任意の \(s, t \in [0,1]\) に対して
\(\lvert t - s\rvert \le \rho(d_t, d_s) \le 2\,\lvert t - s\rvert\)．}

\emph{証明.} 下界は
\(\lvert d_t(p,q) - d_s(p,q)\rvert = \lvert(1+t)-(1+s)\rvert = \lvert t-s\rvert\)
による．上界は，各非順序対で
\(\lvert d_t(u,v) - d_s(u,v)\rvert \le 2\lvert t-s\rvert\)
を確かめて上限をとればよい．\(\{p,q\}\) では差は
\(\lvert t-s\rvert\)．\(r \mapsto \min(1,r)\) と
\(r \mapsto \lvert r - a_\alpha\rvert + 1/(n+1)\) はともに
\(1\)-Lipschitz だから合成 \(r \mapsto s_n(\alpha,r)\) も
\(1\)-Lipschitz であり，\(\{y_\alpha, x_{\alpha,n}\}\) では差は高々
\(\lvert t-s\rvert\)，\(\{x_{\alpha,n}, x_{\alpha,m}\}\) では高々
\(2\lvert t-s\rvert\)．その他の対では値は \(t\) によらず \(1\) で，差は
\(0\)．

\textbf{系.} (C2) より \(\varphi(t) := d_t\) は \(A^\complement\) から
\(\operatorname{Met}(D)\) への写像を定める．(C4) の下界より \(\varphi\)
は単射であり，両側評価より \(\varphi\) は \(A^\complement\) から像
\(\varphi(A^\complement)\)（\(\operatorname{Met}(D)\)
の相対位相を与える）への双リプシッツ同相，すなわち双リプシッツ埋め込みである．

\textbf{(C5) \(\varphi(A^\complement)\) は \(\operatorname{Met}(D)\)
の閉集合である．}

\emph{証明.} \(\operatorname{Met}(D)\) の位相は距離 \(\min(\rho,1)\)
が定めるから，閉包の点は点列の極限として得られる．\(d' \in \operatorname{Met}(D)\)
を \(\varphi(A^\complement)\)
の閉包の点とし，\(t_k \in A^\complement\)（\(k \in \mathbf{N}\)）を
\(\rho(d_{t_k}, d') \to 0\)
となるように取る．\(d' \in \varphi(A^\complement)\) を示す．

\textbf{手順 1（パラメータの回収）.}
\(\lvert d_{t_k}(p,q) - d'(p,q)\rvert \le \rho(d_{t_k}, d') \to 0\) より
\(1 + t_k \to d'(p,q)\)，すなわち \(t_k \to t := d'(p,q) - 1\)
である．各 \(t_k \in [0,1]\) で \([0,1]\) は閉だから \(t \in [0,1]\)
であり，(C1) より \(d_t\) は \(D\)
上の距離として定まる（\(d_t \in \operatorname{Met}(D)\)
はまだ主張しない）．

\textbf{手順 2（極限の同定）.} \(u, v \in D\) を任意に固定する．(C4)
より \(\rho(d_{t_k}, d_t) \le 2\lvert t_k - t\rvert \to 0\) だから
\[\lvert d'(u,v) - d_t(u,v)\rvert \le \rho(d', d_{t_k}) + \rho(d_{t_k}, d_t) \longrightarrow 0 \quad (k \to \infty)\]
であり，左辺は \(k\) によらないから \(d'(u,v) = d_t(u,v)\)．よって
\(d' = d_t\)．

\textbf{手順 3（\(t \in A\) の排除）.} 仮定より
\(d' \in \operatorname{Met}(D)\) だから，手順 2 より
\(d_t = d' \in \operatorname{Met}(D)\) である．\(t \in A\) ならば (C3)
より \(d_t \notin \operatorname{Met}(D)\)
となって矛盾するから，\(t \in A^\complement\) である．

以上より \(d' = d_t = \varphi(t) \in \varphi(A^\complement)\)
となり，\(\varphi(A^\complement)\) は閉である．

パラメータの極限が \(A\)
に落ちる場合を排除できることが閉性の本質であり，それを与えるのが (C3)
である．(C4) の系と (C5) より補題 3.4 が示された．

\hypertarget{ux4e3bux5b9aux7406ux306eux8a3cux660e}{%
\subsection{4.
主定理の証明}\label{ux4e3bux5b9aux7406ux306eux8a3cux660e}}

\textbf{定理 2.1 の証明.} §3.2 より，\(D\) としては §3.3
で構成した具体的な集合を取ってよい．

\begin{itemize}
\tightlist
\item
  \textbf{ステップ 1.} 補題 3.2 の
  \(A \subseteq [0,1]\)（\(\lvert A\rvert = \aleph_1\)，\(A\) は
  \(F_\sigma\) でない）をとる．
\item
  \textbf{ステップ 2.} \(F_\sigma\) と \(G_\delta\)
  の補集合双対により，\(A\) が \([0,1]\) で \(F_\sigma\)
  でないことから，\(A^\complement\) は \([0,1]\) で \(G_\delta\)
  でない．
\item
  \textbf{ステップ 3.} 定理 1.1 の対偶を距離化可能空間 \(X = [0,1]\)
  に適用して，\(A^\complement\) は \([0,1]\)
  の部分空間として完備距離化可能でない．
\item
  \textbf{ステップ 4.} 補題 3.3・補題 3.4
  より，\(\varphi : A^\complement \to \operatorname{Met}(D)\)，\(t \mapsto d_t\)
  は像 \(\varphi(A^\complement)\)
  への同相写像であり，\(\varphi(A^\complement)\) は
  \(\operatorname{Met}(D)\) の閉部分集合である．
\item
  \textbf{ステップ 5.} \(\operatorname{Met}(D)\)
  が完備距離化可能であると仮定すると，定理 1.2 より閉部分集合
  \(\varphi(A^\complement)\)
  も完備距離化可能であり，完備距離化可能性は同相で保たれるから，ステップ
  4 の同相より \(A^\complement\) も完備距離化可能となる．これはステップ
  3 に矛盾する．
\item
  \textbf{ステップ 6.} よって \(\operatorname{Met}(D)\)
  は完備距離化可能でない．
\end{itemize}

\textbf{Remark（濃度条件 \(\lvert A\rvert = \aleph_1\) の使用箇所）.}
補題 3.2 の与える 2 条件の役割は分かれている．\(A\) が \(F_\sigma\)
でないことは証明のステップ 2--3 でのみ用いる．濃度条件
\(\lvert A\rvert = \aleph_1\) は §3.3 の構成の 2 箇所，すなわち \(A\) の
\(\omega_1\) による単射列挙と，濃度計算
\(\lvert D\rvert = 2 + \aleph_1 + \aleph_1 \cdot \aleph_0 = \aleph_1\)
とで用いる．構成がひとたび与えられた後は，(C1)--(C5)
の検証は濃度を用いない．

\hypertarget{ux53c2ux8003ux6587ux732e}{%
\subsection{5. 参考文献}\label{ux53c2ux8003ux6587ux732e}}

\begin{enumerate}
\def\labelenumi{\arabic{enumi}.}
\tightlist
\item
  Y. Ishiki, \emph{An isometric extensor of metrics}, arXiv:2407.03030
  (2024).
\item
  A. S. Kechris, \emph{Classical Descriptive Set Theory}, Graduate Texts
  in Mathematics 156, Springer-Verlag, 1995.
\item
  R. Engelking, \emph{General Topology}, revised ed., Sigma Series in
  Pure Mathematics 6, Heldermann Verlag, 1989.
\item
  S. Willard, \emph{General Topology}, Addison-Wesley, 1970（Dover 版
  2004）.
\item
  J. C. Oxtoby, \emph{Measure and Category}, 2nd ed., Graduate Texts in
  Mathematics 2, Springer-Verlag, 1980.
\end{enumerate}

\end{document}
